\documentclass[a4paper,12pt]{article}
\usepackage[T2A]{fontenc}
\usepackage[utf8]{inputenc}
\usepackage[russian]{babel}
\usepackage{amsmath,amssymb,mathtools}
\renewcommand{\familydefault}{\sfdefault}
\usepackage{icomma}
\usepackage{geometry}
\geometry{left=24mm,right=24mm,top=24mm,bottom=24mm}
\usepackage{microtype}
\usepackage{xcolor}
\usepackage{tikz}
\usetikzlibrary{calc,arrows.meta,positioning,shapes.geometric}
\usepackage{booktabs,tabularx}
\usepackage{enumitem}
\usepackage{parskip}
\usepackage{fancyhdr}
\usepackage{setspace}
\usepackage{hyperref}
\hypersetup{hidelinks}

\definecolor{accent}{HTML}{176B70}
\definecolor{darkaccent}{HTML}{194F55}
\definecolor{textgray}{HTML}{333333}
\color{textgray}
\onehalfspacing
\setlength{\headheight}{16pt}
\pagestyle{fancy}
\fancyhf{}
\fancyhead[L]{\small Ксения Васильченко}
\fancyhead[R]{\small 08.10.26}
\fancyfoot[C]{\small\thepage}
\renewcommand{\headrulewidth}{0.35pt}
\renewcommand{\footrulewidth}{0pt}
\renewcommand{\arraystretch}{1.18}
\setlist[enumerate]{leftmargin=*,itemsep=.55em,topsep=.25em}
\newcommand{\sectionline}[1]{\section*{\color{darkaccent}#1}}
\newcommand{\learn}[1]{\par\noindent$\square$\enspace #1\par}
\newcommand{\keyidea}[1]{\par\noindent\textcolor{darkaccent}{\textbf{Запомните.}}\enspace #1\par}

\begin{document}
\begin{titlepage}
\thispagestyle{empty}
\centering
\vspace*{32mm}
{\color{accent}\large МАТЕМАТИКА \textbullet\ 6 КЛАСС\par}
\vspace{19mm}
{\color{darkaccent}\Huge\bfseries Чек-лист\par}
\vspace{8mm}
{\LARGE Сравнение дробных величин\par}
\vspace{3mm}
{\Large и текстовые задачи на работу\par}
\vspace{21mm}
{\large 8 октября 2026 года\par}
\vfill
{\large Лёвин Артём Александрович\par}
\vspace{2mm}
{\normalsize эксклюзивно для Ксении Васильченко\par}
\vspace{33mm}
\begin{tikzpicture}[remember picture,overlay]
\draw[accent!58,line width=1.25pt]
  ([xshift=25mm,yshift=20mm]current page.south west)
  .. controls ([xshift=79mm,yshift=42mm]current page.south west)
           and ([xshift=-77mm,yshift=4mm]current page.south east)
  .. ([xshift=-25mm,yshift=20mm]current page.south east);
\end{tikzpicture}
\end{titlepage}

\sectionline{1. Дробная часть минуты}
\learn{Вы умеете переводить минуты в секунды и сравнивать полученные значения.}
\learn{Вы умеете сравнивать дробные части минуты без перевода в секунды.}

\keyidea{В одной минуте 60 секунд. Если время равно $a$ минутам, то оно содержит $60a$ секунд.}
\[
\boxed{t_{\text{с}}=60\,t_{\text{мин}}}
\]

\textbf{Пример.} Сравним $\frac13$ минуты и $\frac25$ минуты двумя способами.

\textit{Способ 1. Перевод в секунды:}
\[
\frac13\cdot60=20\ \text{с},\qquad
\frac25\cdot60=24\ \text{с};\qquad 20<24.
\]

\textit{Способ 2. Общий знаменатель:}
\[
\text{НОК}(3,5)=15,\qquad
\frac13=\frac5{15}<\frac6{15}=\frac25.
\]

\keyidea{Оба способа дают один вывод: $\frac13$ минуты $<$ $\frac25$ минуты. При сравнении величин сначала убедитесь, что их единицы измерения одинаковы.}

\textbf{Ещё один пример.} Сравним $\frac{11}{20}$ и $\frac8{15}$ минуты:
\[
\text{НОК}(20,15)=60,\qquad
\frac{11}{20}=\frac{33}{60}>\frac{32}{60}=\frac8{15}.
\]

\sectionline{2. Как сравнивать доли одного целого}
\learn{Вы выделяете в тексте, какую долю общего количества обозначает каждая дробь.}
\learn{Вы понимаете, когда для ответа достаточно сравнить дроби.}

\textbf{Пример.} В парке берёзы составляют $\frac{11}{20}$ всех деревьев, а ясени~--- $\frac{13}{35}$. Каких деревьев больше?
\[
\text{НОК}(20,35)=140,\qquad
\frac{11}{20}=\frac{77}{140}>\frac{52}{140}=\frac{13}{35}.
\]
\textbf{Ответ:} берёз больше.

\keyidea{Если дроби относятся к \emph{одному и тому же целому}, то большая доля соответствует большему количеству. Для сравнения количеств знать размер целого необязательно.}

\sectionline{3. Перевод текста задачи на язык дробей}
\learn{Вы находите долю как отношение части к целому.}
\learn{Вы делите длину на число равных частей, чтобы найти длину одной части.}

\textbf{Меткость.} Никита попал 8 раз из 15 выстрелов, Ярослав~--- 11 раз из 20. Доли попаданий:
\[
\frac8{15}=\frac{32}{60}<\frac{33}{60}=\frac{11}{20}.
\]
\textbf{Вывод:} доля попаданий у Ярослава больше.

\textbf{Равные отрезки.} Брус длиной 5 м распилили на 7 равных частей, другой длиной 6 м~--- на 10 частей. Длины частей:
\[
\frac57\ \text{м}=\frac{50}{70}\ \text{м}>
\frac{42}{70}\ \text{м}=\frac6{10}\ \text{м}.
\]
\textbf{Вывод:} часть первого бруса длиннее.

\keyidea{Доля $=\frac{\text{часть}}{\text{целое}}$; длина равной части $=\frac{\text{общая длина}}{\text{число частей}}$.}

\sectionline{4. Задачи на работу и производительность}
\learn{Вы различаете результат работы, время и производительность.}
\learn{Вы составляете таблицу и вычисляете результат за нужное время.}
\learn{Вы сравниваете результаты даже тогда, когда размер всего заказа неизвестен.}

Обозначим результат работы буквой $R$, время~--- $t$, производительность (результат за единицу времени)~--- $p$.
\[
\boxed{p=\frac Rt}\qquad
\boxed{R=pt}\qquad
\boxed{t=\frac Rp},
\quad t>0,\ p>0.
\]

\textbf{Пример.} Мастер выполняет один заказ за 6 ч, ученик~--- тот же заказ за 8 ч. Кто изготовит больше деталей: мастер за 5 ч или ученик за 7 ч?

Пусть заказ содержит $n$ деталей, где $n>0$. Тогда:
\[
\begin{array}{@{}lccc@{}}
\toprule
&\text{Весь заказ}&\text{Время}&\text{Производительность}\\
\midrule
\text{Мастер}&n&6\ \text{ч}&\frac n6\\
\text{Ученик}&n&8\ \text{ч}&\frac n8\\
\bottomrule
\end{array}
\]

За указанные промежутки времени результаты равны:
\[
R_{\text{м}}=\frac n6\cdot5=\frac{5n}6=\frac{20n}{24},\qquad
R_{\text{у}}=\frac n8\cdot7=\frac{7n}8=\frac{21n}{24}.
\]
При $n>0$ имеем $\frac{21n}{24}>\frac{20n}{24}$. \textbf{Ответ:} ученик за 7 часов изготовит больше деталей.

\keyidea{Если объём заказа неизвестен, обозначьте его $n>0$. Когда удобнее сравнивать \emph{доли выполненной работы}, весь заказ можно принять за $1$.}

\textbf{Случай равенства.} Первый работник делает заказ за 5 ч, второй~--- за 10 ч. Сравним их результаты за 3 и 6 ч:
\[
\frac n5\cdot3=\frac{3n}5,\qquad
\frac n{10}\cdot6=\frac{3n}5.
\]
\textbf{Вывод:} за эти промежутки времени работники выполнят одинаковую часть заказа.

\textbf{Самопроверка:} при постоянной производительности $p$ удвоение времени удваивает результат работы $R$.

\clearpage
\thispagestyle{plain}
\sectionline{Алгоритм: как сравнить результаты работы}
\vspace{12mm}
\begin{center}
\begin{tikzpicture}[
  >=Stealth,
  every node/.style={align=center,font=\normalsize},
  stepbox/.style={draw=accent,rounded corners=3pt,text width=11.8cm,minimum height=1.48cm,inner sep=7pt},
  terminal/.style={draw=darkaccent,ellipse,minimum width=5cm,minimum height=1.15cm},
  arrow/.style={->,thick,draw=darkaccent},
  node distance=7mm
]
\node[terminal] (a) {Начало};
\node[stepbox,below=of a] (b) {Определите, кто выполняет работу, за какое время и какой заказ};
\node[stepbox,below=of b] (c) {Составьте таблицу: результат $R$, время $t$, производительность $p$};
\node[stepbox,below=of c] (d) {Для одинакового заказа примите $R=n>0$ или $R=1$};
\node[stepbox,below=of d] (e) {Найдите производительность каждого: $p=R/t$};
\node[stepbox,below=of e] (f) {За время из вопроса найдите результат: $R=pt$};
\node[stepbox,below=of f] (g) {Приведите дроби к общему знаменателю и сравните};
\node[terminal,below=of g] (h) {Запишите ответ};
\foreach \x/\y in {a/b,b/c,c/d,d/e,e/f,f/g,g/h}{\draw[arrow] (\x.south)--(\y.north);}
\end{tikzpicture}
\end{center}
\vspace{8mm}
\noindent\textit{Условие применения:} сравниваем работу при постоянной производительности исполнителей.

\clearpage
\sectionline{5. Самостоятельная тренировка \textnormal{(Путь А)}}
\textbf{10 заданий.} Выполните решение письменно. В текстовых задачах укажите, какие дроби Вы сравниваете, и запишите вывод словами.

\begin{enumerate}[itemsep=1.05em]
\item Сколько секунд составляют $\frac34$ минуты?
\item Сравните $\frac5{12}$ минуты и $\frac25$ минуты, выразив оба промежутка в секундах.
\item Сравните $\frac7{18}$ минуты и $\frac5{12}$ минуты, приведя дроби к наименьшему общему знаменателю.
\item В саду яблони составляют $\frac7{16}$ всех деревьев, а груши~--- $\frac5{12}$. Каких деревьев больше?
\item Один участник попал в цель 9 раз из 16 выстрелов, другой~--- 11 раз из 20. У кого выше доля попаданий?
\item Верёвку длиной 7 м разделили на 9 равных частей, а верёвку длиной 5 м~--- на 7 равных частей. Какая получившаяся часть длиннее?
\item Мастер выполняет весь заказ за 5 часов. Какую долю заказа он выполнит за 3 часа при постоянной производительности?
\item Первый работник завершает заказ за 9 часов, второй такой же заказ~--- за 12 часов. Кто выполнит большую часть своего заказа: первый за 4 часа или второй за 6 часов?
\item Первый мастер изготавливает заказ за 5 часов, второй такой же заказ~--- за 10 часов. Сравните результаты первого за 3 часа и второго за 6 часов.
\item Первый насос наполняет бак за 18 часов, второй такой же бак~--- за 15 часов. Какой насос наполнит большую часть бака: первый за 5 часов или второй за 4 часа?
\end{enumerate}

\clearpage
\sectionline{6. Ответы для самопроверки}
\begin{tabularx}{\linewidth}{@{}p{13mm}X@{}}
\toprule
\textbf{№}&\textbf{Ответ и ключевая проверка}\\
\midrule
1&$45$ секунд.\\
2&$\frac5{12}$ минуты больше: $25$ с $>24$ с.\\
3&$\frac7{18}<\frac5{12}$, так как $\frac{14}{36}<\frac{15}{36}$.\\
4&Яблонь больше: $\frac{21}{48}>\frac{20}{48}$.\\
5&У первого участника: $\frac{45}{80}>\frac{44}{80}$.\\
6&Часть первой верёвки длиннее: $\frac{49}{63}>\frac{45}{63}$.\\
7&$\frac35$ заказа.\\
8&Второй работник: $\frac49<\frac12$.\\
9&Одинаковые результаты: $\frac{3n}{5}=\frac{6n}{10}$ при $n>0$.\\
10&Первый насос: $\frac5{18}>\frac4{15}$, так как $\frac{25}{90}>\frac{24}{90}$.\\
\bottomrule
\end{tabularx}

\vspace{9mm}
\keyidea{Проверяя себя, смотрите на смысл ответа: сравниваются доли одного целого или результаты работы за указанные промежутки времени.}
\end{document}
